% Basic Commitment Scheme - Extracted from snargsbook.tex
% Chapter 17: Basic commitment scheme
% Reference: "Building Cryptographic Proofs from Hash Functions" (https://snargsbook.org/)

\documentclass{article}
\usepackage{amsmath,amssymb,amsthm}
\usepackage{paralist}
\usepackage{hyperref}

% Basic formatting
\newtheorem{lemma}{Lemma}
\newtheorem{claim}{Claim}
\newtheorem{remark}{Remark}

% Common math notation
\newcommand{\Naturals}{\mathbb{N}}
\newcommand{\Bits}{\{0,1\}}
\newcommand{\DefineEqual}{\coloneqq}
\newcommand{\Negate}{\overline}
\newcommand{\Cardinality}[1]{|#1|}
\newcommand{\AbsValue}[1]{|#1|}
\newcommand{\Expectation}{\mathbb{E}}
\newcommand{\Pr}{\mathsf{Pr}}
\newcommand{\StatisticalDistance}[2]{\Delta(#1, #2)}
\newcommand{\TextAndInMath}{\quad\text{and}\quad}
\newcommand{\ConditionedOn}{\mid}
\newcommand{\DoQuote}[1]{``#1''}
\newcommand{\EquationComment}[1]{\qquad\text{(#1)}}
\newcommand{\IndentedBullet}{\quad\bullet\ }
\newcommand{\concat}{\,\|\,}

% Experiment notation
\newcommand{\GivenExperiment}{\;\middle|\;}
\newcommand{\StateExperiment}[1]{\begin{array}{l}#1\end{array}}
\newcommand{\Event}{\mathsf{E}}

% Random Oracle notation
\newcommand{\ROFunction}{f}
\newcommand{\RODistribution}[1]{\mathcal{F}_{#1}}
\newcommand{\ROOutputSize}{\lambda}
\newcommand{\ROQueryBound}{q}
\newcommand{\ROTrace}{\mathsf{T}}
\newcommand{\ROAdvState}{\mathsf{st}}
\newcommand{\ROQuery}{\mathsf{x}}
\newcommand{\ROAnswer}{\mathsf{y}}
\newcommand{\ROOutputAndTrace}[4]{(#3, #4) \gets #2^{#1}}
\newcommand{\ROInputOutputAndTrace}[5]{(#4, #5) \gets #2^{#1}(#3)}

% Collision resistance expressions
\newcommand{\ROCRExpression}[2]{\frac{1}{2} \cdot \frac{(#2-1) \cdot #2}{2^{#1}}}
\newcommand{\ROPExpression}[1]{\frac{1}{2^{#1}}}
\newcommand{\ROOWEntropyExpression}[1]{\frac{\ROQueryBound}{2^{#1}}}

% Protocol notation
\newcommand{\ProtoSalt}{\tau}
\newcommand{\ProtoSoundnessError}{\varepsilon}
\newcommand{\ProtoKnowledgeError}{\kappa}
\newcommand{\ProtoAdversary}{\mathcal{A}}

% Basic Commitment Scheme notation
\newcommand{\CMSymbol}{\mathsf{CM}}
\newcommand{\CMSubscript}{\scriptscriptstyle\CMSymbol}
\newcommand{\CMConstructor}[3]{\CMSymbol[#1,#2,#3]}
\newcommand{\CMCommit}{\CMSymbol.\mathsf{Commit}}
\newcommand{\CMCheck}{\CMSymbol.\mathsf{Check}}
\newcommand{\CMExtractor}{\CMSymbol.\mathsf{Extract}}
\newcommand{\CMMultiExtractor}{\CMSymbol.\mathsf{MultiExtract}}
\newcommand{\CMSimulator}{\CMSymbol.\mathsf{Simulate}}
\newcommand{\CMMessageLength}{\ell}
\newcommand{\CMSaltSize}{s}
\newcommand{\CMCommitment}{\mathsf{cm}}
\newcommand{\CMMessage}{\mathsf{m}}
\newcommand{\CMSaltString}{\ProtoSalt}
\newcommand{\CMBindingError}{\ProtoSoundnessError_{\CMSubscript}}
\newcommand{\CMExtractionError}{\ProtoKnowledgeError_{\CMSubscript}}
\newcommand{\CMHidingError}{z_{\CMSubscript}}
\newcommand{\CMNumCommitments}{n}
\newcommand{\CMAdversary}{\ProtoAdversary}

% Error bound expressions
\newcommand{\CMBindingExpression}[2]{\frac{1}{2} \cdot \frac{#2^2}{2^{#1}}}
\newcommand{\CMExtractabilityExpression}[2]{\CMBindingExpression{#1}{#2}}
\newcommand{\CMMultiExtractabilityExpression}[3]{\frac{1}{2} \cdot \frac{#2^2}{2^{#1}} + \frac{#3 \cdot #2}{2^{#1}}}
\newcommand{\CMHidingExpressionFinite}[2]{\frac{#2}{2^{#1}}}
\newcommand{\CMHidingExpressionInfinite}[1]{3 \sqrt{\ROOutputSize \CMMessageLength #1} \cdot 2^{-\frac{#1 - \ROOutputSize}{2}}}

% Random variable
\newcommand{\RandomVariableZ}{Z}
\newcommand{\RandomVariableX}{X}

% Cross-references (simplified)
\newcommand{\Cref}[1]{Section~\ref{#1}}
\newcommand{\FigureFolder}{figures}

\title{Basic Commitment Scheme}
\author{Extracted from ``Building Cryptographic Proofs from Hash Functions''}
\date{}

\begin{document}
\maketitle

\begin{abstract}
This document contains the basic commitment scheme chapter from the snargsbook,
describing a commitment scheme $\CMSymbol = (\CMCommit, \CMCheck)$ in the random
oracle model. The scheme satisfies three key properties: binding, extractability,
and hiding.
\end{abstract}

\tableofcontents

\section{Introduction}
\label{chapter:basic-commitment}

A (non-interactive) commitment scheme enables a party to commit to a message $\CMMessage$ to obtain a commitment $\CMCommitment$ in such a way that:
\begin{inparaenum}[(a)]
  \item the committing party can subsequently ``open'' the commitment $\CMCommitment$ to the message $\CMMessage$;
  \item it is intractable to open $\CMCommitment$ to a different message $\CMMessage'$; and
  \item the commitment $\CMCommitment$ reveals no information about the message $\CMMessage$.
\end{inparaenum}
Commitment schemes are a fundamental cryptographic primitive, and in particular they play a key role in the construction of succinct arguments in the random oracle model.

\subsection{The Basic Construction}

The basic commitment scheme is a tuple of algorithms
\begin{equation*}
\CMSymbol = (\CMCommit,\CMCheck)
\end{equation*}
and has several parameters: an output size $\ROOutputSize \in \Naturals$ of the random oracle, a message length $\CMMessageLength \in \Naturals$, and a salt size $\CMSaltSize \in \Naturals$. We use the notation $\CMConstructor{\ROOutputSize}{\CMMessageLength}{\CMSaltSize}$ when we wish to make these parameters explicit. The algorithms receive query access to a random oracle $\ROFunction \in \RODistribution{\ROOutputSize}$ and work as follows.

\begin{itemize}
  \item \emph{Commit.}
  The algorithm $\CMCommit$ receives as input a message $\CMMessage \in \Bits^{\CMMessageLength}$, samples a random salt $\CMSaltString \in \Bits^{\CMSaltSize}$, queries the random oracle $\ROFunction$ to compute the commitment $\CMCommitment \DefineEqual \ROFunction(\CMMessage,\CMSaltString) \in \Bits^{\ROOutputSize}$, and outputs the pair $(\CMCommitment,\CMSaltString)$.

  \item \emph{Check.}
  The algorithm $\CMCheck$ receives as input a commitment $\CMCommitment$, message $\CMMessage \in \Bits^{\CMMessageLength}$, and salt $\CMSaltString \in \Bits^{\CMSaltSize}$, and queries the random oracle to check that $\CMCommitment = \ROFunction(\CMMessage,\CMSaltString)$.
\end{itemize}

\subsection{Security Properties Overview}

We prove several properties about the basic commitment scheme $\CMSymbol$:
\begin{itemize}
  \item \textbf{Binding} (Section~\ref{section:basic-commitment-binding}): An adversary cannot find two distinct messages that open the same commitment. The binding error is at most $\CMBindingExpression{\ROOutputSize}{\ROQueryBound}$ for $\ROQueryBound$-query adversaries.

  \item \textbf{Extractability} (Section~\ref{section:basic-commitment-extractability}): If an adversary outputs a commitment that it subsequently opens, then the adversary ``knew'' the opening at commitment time. The extraction error equals the binding error.

  \item \textbf{Hiding} (Section~\ref{section:basic-commitment-hiding}): A commitment reveals no information about the committed message. The hiding error is at most $\CMHidingExpressionFinite{\CMSaltSize}{\ROQueryBound}$ for bounded adversaries.
\end{itemize}

\section{Binding}
\label{section:basic-commitment-binding}

We prove that $\CMSymbol$ is \emph{binding}: an adversary cannot find two distinct messages that open the same commitment.

\begin{lemma}[$\CMSymbol$ is binding]
\label{lemma:cm-binding}
Let $\CMSymbol \DefineEqual \CMConstructor{\ROOutputSize}{\CMMessageLength}{\CMSaltSize}$. For every query bound $\ROQueryBound \in \Naturals$ and $\ROQueryBound$-query algorithm $\CMAdversary$, if $\ROQueryBound \geq 2$ then
\begin{equation*}
\Pr\left[
\begin{array}{l}
\CMMessage_{0} \neq \CMMessage_{1} \\
\land\;\CMCheck^{\ROFunction}(\CMCommitment,\CMMessage_{0},\CMSaltString_{0})=1 \\
\land\;\CMCheck^{\ROFunction}(\CMCommitment,\CMMessage_{1},\CMSaltString_{1})=1
\end{array}
\GivenExperiment
\StateExperiment{
\ROFunction \gets \RODistribution{\ROOutputSize} \\
(\CMCommitment,\CMMessage_{0},\CMSaltString_{0},\CMMessage_{1},\CMSaltString_{1}) \gets \CMAdversary^{\ROFunction}
}
\right]
\leq
\CMBindingExpression{\ROOutputSize}{\ROQueryBound}
\end{equation*}
\end{lemma}

The proof relies on collision resistance of the random oracle and on unpredictability for edge cases.

\section{Extractability}
\label{section:basic-commitment-extractability}

We prove that $\CMSymbol$ is \emph{extractable}: if an adversary outputs a commitment that it subsequently opens, then the adversary ``knew'' the opening at commitment time.

\subsection{Extraction for One Commitment}

The extractability property is formalized via an efficient algorithm $\CMExtractor$ called an \emph{extractor}.

\begin{lemma}[$\CMSymbol$ is extractable]
\label{lemma:cm-extractability}
Let $\CMSymbol \DefineEqual \CMConstructor{\ROOutputSize}{\CMMessageLength}{\CMSaltSize}$. There exists a polynomial-time deterministic algorithm $\CMExtractor$ such that for every query bound $\ROQueryBound \in \Naturals$ and $\ROQueryBound$-query algorithm $\CMAdversary$, if $\ROQueryBound \geq 3$ then
\begin{equation*}
\Pr\left[
\begin{array}{l}
\CMCheck^{\ROFunction}(\CMCommitment,\CMMessage,\CMSaltString)=1 \\
\land\;(\CMMessage',\CMSaltString') \neq (\CMMessage,\CMSaltString)
\end{array}
\GivenExperiment
\StateExperiment{
\ROFunction \gets \RODistribution{\ROOutputSize} \\
\ROOutputAndTrace{\ROFunction}{\CMAdversary}{\ROTrace}{(\CMCommitment,\ROAdvState)} \\
(\CMMessage',\CMSaltString') \gets \CMExtractor(\CMCommitment,\ROTrace) \\
(\CMMessage,\CMSaltString) \gets \CMAdversary^{\ROFunction}(\ROAdvState)
}
\right]
\leq
\CMExtractabilityExpression{\ROOutputSize}{\ROQueryBound}
\end{equation*}
Moreover, $\CMExtractor$ runs in time $O\big(\ROQueryBound \cdot (\CMMessageLength + \CMSaltSize + \ROOutputSize)\big)$.
\end{lemma}

The extractor $\CMExtractor$ works as follows: Search the trace $\ROTrace$ for a query-answer pair of the form $((\CMMessage',\CMSaltString'),\CMCommitment)$. If found, output $(\CMMessage',\CMSaltString')$; otherwise abort.

\subsection{Extraction for Multiple Commitments}

Extraction extends to adversaries that output \emph{multiple} commitments without multiplicative loss in the number of commitments.

\begin{lemma}[$\CMSymbol$ is multi-extractable]
\label{lemma:cm-multi-extractability}
Let $\CMSymbol \DefineEqual \CMConstructor{\ROOutputSize}{\CMMessageLength}{\CMSaltSize}$. There exists a polynomial-time deterministic \textbf{stateful} algorithm $\CMMultiExtractor$ such that for every query bound $\ROQueryBound \in \Naturals$, $\ROQueryBound$-query algorithm $\CMAdversary$, and number of commitments $\CMNumCommitments \in \Naturals$, if $\ROQueryBound \geq 2\CMNumCommitments$ then the extraction error is at most
\begin{equation*}
\CMMultiExtractabilityExpression{\ROOutputSize}{\ROQueryBound}{\CMNumCommitments}
\end{equation*}
\end{lemma}

The dominant error (collisions) is global rather than per-commitment, which is why there is no multiplicative loss.

\subsection{Extractability Implies Binding}

Extractability is stronger than binding: if $\CMSymbol$ satisfies extractability with error $\CMExtractionError$ then $\CMSymbol$ satisfies binding with error $\CMBindingError \leq 2 \cdot \CMExtractionError$.

\section{Hiding}
\label{section:basic-commitment-hiding}

We prove that $\CMSymbol$ is \emph{hiding}: a commitment $\CMCommitment$ computed as $\ROFunction(\CMMessage,\CMSaltString)$ for a random $\CMSaltString \in \Bits^{\CMSaltSize}$ reveals no information about $\CMMessage$, up to a small statistical ``leakage'' error.

\begin{lemma}[$\CMSymbol$ is hiding]
\label{lemma:cm-hiding}
Let $\CMSymbol \DefineEqual \CMConstructor{\ROOutputSize}{\CMMessageLength}{\CMSaltSize}$. There exists a polynomial-time probabilistic algorithm $\CMSimulator$ such that for every query bound $\ROQueryBound \in \Naturals$ and $\ROQueryBound$-query algorithm $\CMAdversary$, the following two distributions are $\CMHidingError(\ROOutputSize,\CMMessageLength,\CMSaltSize,\ROQueryBound)$-close in statistical distance:
\begin{equation*}
\left\{
\CMAdversary^{\ROFunction}(\ROAdvState,\CMCommitment)
\GivenExperiment
\StateExperiment{
\ROFunction \gets \RODistribution{\ROOutputSize} \\
(\CMMessage,\ROAdvState) \gets \CMAdversary^{\ROFunction} \\
(\CMCommitment,\CMSaltString) \gets \CMCommit^{\ROFunction}(\CMMessage)
}
\right\}
\TextAndInMath
\left\{
\CMAdversary^{\ROFunction}(\ROAdvState,\CMCommitment)
\GivenExperiment
\StateExperiment{
\ROFunction \gets \RODistribution{\ROOutputSize} \\
(\CMMessage,\ROAdvState) \gets \CMAdversary^{\ROFunction} \\
\CMCommitment \gets \CMSimulator^{\ROFunction}
}
\right\}
\end{equation*}
where
\begin{equation*}
\CMHidingError(\ROOutputSize,\CMMessageLength,\CMSaltSize,\ROQueryBound)
\leq
\begin{cases}
\CMHidingExpressionFinite{\CMSaltSize}{\ROQueryBound}
& \text{ if $\ROQueryBound<\infty$} \\
\CMHidingExpressionInfinite{\CMSaltSize}
& \text{ if $\ROQueryBound=\infty$}
\end{cases}
\end{equation*}
Moreover, $\CMSimulator^{\ROFunction}$ outputs a random string in $\Bits^{\ROOutputSize}$.
\end{lemma}

\section{Summary of Error Bounds}

\begin{center}
\begin{tabular}{|l|l|}
\hline
\textbf{Property} & \textbf{Error Bound} \\
\hline
Binding & $\CMBindingExpression{\ROOutputSize}{\ROQueryBound}$ \\
Extractability & $\CMExtractabilityExpression{\ROOutputSize}{\ROQueryBound}$ \\
Multi-Extractability & $\CMMultiExtractabilityExpression{\ROOutputSize}{\ROQueryBound}{\CMNumCommitments}$ \\
Hiding (bounded) & $\CMHidingExpressionFinite{\CMSaltSize}{\ROQueryBound}$ \\
Hiding (unbounded) & $\CMHidingExpressionInfinite{\CMSaltSize}$ \\
\hline
\end{tabular}
\end{center}

\section*{References}

This chapter is based on Chapter 17 of ``Building Cryptographic Proofs from Hash Functions'' available at \url{https://snargsbook.org/}.

\end{document}
